\question[4]
Beim One-Time-Pad wird ein Klartext mit einem zufällig erzeugten 
Schlüssel verschlüsselt, der genauso lang wie der Klartext ist und nur 
ein einziges Mal verwendet wird.

\begin{enumerate}
    \item Erkläre kurz, wie die Ver- und Entschlüsselung beim 
    One-Time-Pad grundsätzlich abläuft.
    \item Begründe, warum das One-Time-Pad als absolut sicheres 
    Verschlüsselungsverfahren gilt, sofern der Schlüssel wirklich 
    zufällig, geheim und nur einmal verwendet wird.
    \item Nenne einen Grund, warum das One-Time-Pad trotz seiner 
    absoluten Sicherheit in der Praxis kaum eingesetzt wird.
\end{enumerate}

\begin{solutionbox}{10cm}

\textbf{1.} Klartext und Schlüssel bestehen aus Zeichen gleicher Länge. 
Jedes Klartextzeichen wird mit dem an derselben Position stehenden 
Schlüsselzeichen verknüpft, wodurch das Geheimtextzeichen 
entsteht. Zur Entschlüsselung wird derselbe Schlüssel mit der 
Umkehroperation angewendet.

\textbf{2.} Da der Schlüssel echt zufällig, genauso lang wie der 
Klartext und nur ein einziges Mal verwendet wird, ist für jedes 
mögliche Geheimtextzeichen jedes Klartextzeichen gleich wahrscheinlich. 
Ein Angreifer, der nur den Geheimtext kennt, kann somit rein rechnerisch 
keinerlei Information über den Klartext gewinnen, egal wie viel 
Rechenleistung ihm zur Verfügung steht -- selbst Ausprobieren aller 
Schlüssel liefert jeden denkbaren Klartext gleicher Länge als mögliches 
Ergebnis. 

\textbf{3.} Der Schlüssel muss mindestens so lang wie der Klartext sein 
und darf nur ein einziges Mal benutzt werden. Für einen sicheren 
Nachrichtenaustausch müsste daher für jede Nachricht ein neuer, echt 
zufälliger Schlüssel erzeugt und auf einem sicheren Weg zwischen Sender 
und Empfänger ausgetauscht werden, was in der Praxis (großer 
Schlüsselaufwand, sichere Verteilung) meist unpraktikabel ist. 


\end{solutionbox}